\section{Band selection: which method where, from the bounds}
\label{sec:join}

The bounds of Sections~\ref{sec:bnd} and~\ref{sec:bnd2} are enough to decide, for every
relative offset of every cell shape, which method computes that offset and with how many
terms, before any value is computed. This section states the rule, gives the cost model it
minimises, and tabulates the band edges the bounds produce. Nothing here is fitted to a
measurement; measurement enters only through the cost model and through the term cap
$L_{\max}$, and both are labelled where they do.

Two sets of band-edge numbers appear below and it is worth separating them once.
Section~\ref{sec:join:bands} tabulates the output of \code{bands.jl}, which evaluates the
\emph{first-round} bounds \eqref{eq:bnd:rem} and \eqref{eq:bnd:sub} at
$\mathrm{tol} = 10^{-13}$ with a cap $L_{\max} = 64$; it is what established, before the
library existed, that $L$ is governed by $\rho$ and not by $k\mathbf{R}$ and that exactly
twelve offsets per octant need subdivision. Section~\ref{sec:join:ship} then gives the
edges the \emph{shipped} selector produces, from Theorem~A and Theorem~B of
Section~\ref{sec:bnd2} at $\mathrm{tol} = 10^{-14}$ with $L_{\max} = 56$. The two agree on
every qualitative conclusion and differ in $L$ by $-17\,\%$ to $+17\,\%$.

\subsection{The three methods and their costs}
\label{sec:join:cost}

\begin{description}[nosep,leftmargin=1.4em]
\item[\textbf{W} --- whole box.] Equation~\eqref{eq:exp:eval} with $L = L_{\mathrm{wb}}$
from Theorem~\ref{thm:bnd:rem}. Measured cost $391$, $1250$, $2731$, $4972$~ns at
$L = 10, 20, 30, 40$, i.e.
\begin{equation}
c_{\mathrm{W}}(L) \approx 3.1\,L^{2} + 80 \ \text{ns},
\label{eq:join:cw}
\end{equation}
about $3$~ns per $(l,m)$ slot at roughly $1.7$ complex multiplies each, with no
allocation.
\item[\textbf{O} --- octant split.] Equation~\eqref{eq:exp:subeval} over eight sub-boxes
with $L_j$ from \eqref{eq:bnd:sub}. Seven complex multiplies per $(l,m)$ instead of
$1.7$, so at the same measured $3$~ns per $(l,m)$ per $1.7$ multiplies,
\begin{equation}
c_{\mathrm{O}}(\{L_j\}) \approx 12.2 \sum_{j=1}^{8}(L_j+1)^{2} + 640 \ \text{ns}.
\label{eq:join:co}
\end{equation}
\item[\textbf{K} --- $k$-series.] The moment series of the companion document on the
thirty-six face pairs, with $N$ from Theorem~\ref{thm:bnd:kser}. Cost is dominated by the
moments and is three to four orders larger per offset: $32$~ms at two cells,
$\lambda/32$, in \code{Float64}, and about $0.2$~s per offset in $128$-bit arithmetic,
which is what the amplification $\Lambda$ forces on the slender cell. The moments are
frequency independent, so on a fixed shape this is a one-off per offset.
\end{description}

For reference, the rule being replaced --- thirty-six face-pair integrals by the fixed
Gauss--Legendre schedule of \code{quadOrd} --- costs $13.9$~ms per offset at two cells,
$\lambda/32$, and $0.55$~ms at thirty-two cells, and $23.3$~s for a $32^3$ block.

\emph{Caveat on all timings.} Every microsecond figure quoted here was measured while
several other processes were running on the same machine. Relative comparisons within one
run are sound; absolute per-offset costs must be re-measured on a quiet machine before
they are used for anything but ordering the methods.

Since $c_{\mathrm{W}}(L) < c_{\mathrm{O}}(\{L_j\})$ whenever
$L \lesssim 2\max_j L_j$, and $c_{\mathrm{K}}$ exceeds both by three orders, the cost
ordering is $\mathrm{W} < \mathrm{O} \ll \mathrm{K}$ at every offset where more than one
of them applies. The selection rule is therefore a cascade.

\subsection{The rule}
\label{sec:join:rule}

\begin{equation}
\boxed{
\begin{aligned}
&\text{Take the cheapest method whose \emph{proven} bound reaches }
\mathrm{tol}\cdot\mathrm{est}(\mathbf{R}) \text{ within the term cap } L_{\max}:\\[2pt]
&\quad\text{(1) } L_{\mathrm{wb}}(\mathbf{R}) \le L_{\max}
 \;\Longrightarrow\; \mathrm{W} \text{ at order } L_{\mathrm{wb}};\\
&\quad\text{(2) else } \max_j L_j(\mathbf{R}) \le L_{\max}
 \;\Longrightarrow\; \mathrm{O} \text{ with per-sub-box orders } L_j;\\
&\quad\text{(3) else } \mathrm{K} \text{ at order } N \text{ from
 Theorem~\ref{thm:bnd:kser},}
\end{aligned}}
\label{eq:join:rule}
\end{equation}
with $N$ per shell from Theorem~\ref{thm:bnd:n} in cases (1) and (2)
(Remark~\ref{rem:bnd:adapt}). In the shipped library
$L_{\mathrm{wb}}$ comes from Theorem~A (Theorem~\ref{thm:bnd2:reg}) and $L_j$ from
Theorem~B (Theorem~\ref{thm:bnd2:sub}), with

\begin{center}
\begin{tabular}{@{}ll@{}}
\toprule
$\mathrm{tol} = 10^{-14}$ & \code{TOL} in \code{farfield.jl}\\
$L_{\max} = 56$ & \code{LMAX}\\
$N_{\mathrm{cap}} = 40$ & \code{NCAP}; the call errors above it\\
\bottomrule
\end{tabular}
\end{center}

\noindent
\emph{Why $\mathrm{tol} = 10^{-14}$ and not $10^{-13}$.} The criterion is
$\mathrm{bound} \le \mathrm{tol}\cdot\mathrm{est}$, and $\mathrm{est}$ over-states
$\max_{ab}|T_{ab}|$ by a measured factor $1.10$ to $6.34$ on the cubes
(Section~\ref{sec:unc}), so what the rule certifies is a relative error of
$\mathrm{tol}\cdot\mathrm{est}/\max|T|$. With the first-round bound, loose by a median
$1.5\times10^{5}$, that slack was hidden; with Theorem~A, loose by a median $49$, it is
not. Measured on the same references, $\mathrm{tol} = 10^{-13}$ delivers
$1.6\times10^{-13}$ per entry and $\mathrm{tol} = 10^{-14}$ delivers $9.3\times10^{-15}$,
for $+0.7\,\%$ in terms. The tolerance was tightened because the bound was sharpened.

$L_{\max}$ is the one place where measurement, not proof, enters the rule, and the
justification is threefold, each part a measurement: the exact geometry table costs
$14$--$18$~s per cell shape at $L \approx 50$ and grows steeply beyond; the whole-box
form was measured to reach $6.6\times10^{-14}$ at $L = 94$ at the worst offset and
\emph{not} to improve to $10^{-15}$ by $L = 100$, so a whole-box sum of order
$\sim\!100$ has a \code{Float64} floor above the target; and the octant form reaches
$4.4\times10^{-15}$ with $L_j \le 52$ at exactly those offsets. A cap of $56$ sits above
every $L_j$ the sub-box bound requires on the shapes Gila uses and below the region where
the whole-box route stops paying.

\subsection{Band edges for the four shapes, from the first-round bounds}
\label{sec:join:bands}

$L_{\mathrm{wb}}$ is the whole-box cap from \eqref{eq:bnd:rem}, even $l$ only;
$L_{\mathrm{oct}} = \max_j L_j$ from \eqref{eq:bnd:sub} over the eight octant sub-boxes;
costs from \eqref{eq:join:cw} and \eqref{eq:join:co} in ns; $f = 1$. Method by
\eqref{eq:join:rule} at $\mathrm{tol} = 10^{-13}$ with $L_{\max} = 64$. The shipped
selector's own numbers, from Theorems A and B at $10^{-14}$, are in
Section~\ref{sec:join:ship}.

\begin{center}\footnotesize
\begin{tabular}{@{}llrrrrrrrl@{}}
\toprule
shape & direction & $n$ & $\rho$ & $|k\mathbf{R}|$
& $L_{\mathrm{wb}}$ & $c_{\mathrm{W}}$ & $L_{\mathrm{oct}}$ & $c_{\mathrm{O}}$ & method\\
\midrule
$(1/32)^3$ & axis      & 2  & 0.8660 & 0.393  & 174 & 93936 & 55 & 203648 & O\\
$(1/32)^3$ & axis      & 3  & 0.5774 & 0.589  & 50  & 7830  & 32 & 84283  & W\\
$(1/32)^3$ & axis      & 4  & 0.4330 & 0.785  & 34  & 3664  & 25 & 57248  & W\\
$(1/32)^3$ & axis      & 6  & 0.2887 & 1.178  & 24  & 1866  & 19 & 35971  & W\\
$(1/32)^3$ & axis      & 8  & 0.2165 & 1.571  & 20  & 1320  & 16 & 27236  & W\\
$(1/32)^3$ & axis      & 16 & 0.1083 & 3.142  & 14  & 688   & 12 & 17134  & W\\
$(1/32)^3$ & axis      & 32 & 0.0541 & 6.283  & 12  & 526   & 10 & 12450  & W\\
$(1/32)^3$ & axis      & 64 & 0.0271 & 12.566 & 10  & 390   & 9  & 10400  & W\\
$(1/32)^3$ & face diag & 2  & 0.6124 & 0.555  & 56  & 9802  & 39 & 95971  & W\\
$(1/32)^3$ & face diag & 4  & 0.3062 & 1.111  & 24  & 1866  & 20 & 36923  & W\\
$(1/32)^3$ & face diag & 8  & 0.1531 & 2.221  & 16  & 874   & 14 & 21892  & W\\
$(1/32)^3$ & body diag & 2  & 0.5000 & 0.680  & 40  & 5040  & 32 & 67594  & W\\
$(1/32)^3$ & body diag & 4  & 0.2500 & 1.360  & 22  & 1580  & 18 & 31006  & W\\
$(1/32)^3$ & body diag & 8  & 0.1250 & 2.721  & 14  & 688   & 13 & 19440  & W\\
\midrule
$(1/8)^3$  & axis      & 2  & 0.8660 & 1.571  & 176 & 96106 & 55 & 206820 & O\\
$(1/8)^3$  & axis      & 3  & 0.5774 & 2.356  & 52  & 8462  & 32 & 86772  & W\\
$(1/8)^3$  & axis      & 4  & 0.4330 & 3.142  & 36  & 4098  & 25 & 57248  & W\\
$(1/8)^3$  & axis      & 8  & 0.2165 & 6.283  & 22  & 1580  & 17 & 30554  & W\\
$(1/8)^3$  & axis      & 32 & 0.0541 & 25.13  & 14  & 688   & 12 & 17134  & W\\
$(1/8)^3$  & face diag & 2  & 0.6124 & 2.221  & 58  & 10508 & 39 & 100094 & W\\
$(1/8)^3$  & body diag & 2  & 0.5000 & 2.721  & 42  & 5548  & 33 & 72766  & W\\
\midrule
$(1/4)^3$  & axis      & 2  & 0.8660 & 3.142  & 182 & 102764& 56 & 215604 & O\\
$(1/4)^3$  & axis      & 3  & 0.5774 & 4.712  & 54  & 9120  & 33 & 92628  & W\\
$(1/4)^3$  & axis      & 4  & 0.4330 & 6.283  & 38  & 4556  & 26 & 64324  & W\\
$(1/4)^3$  & axis      & 8  & 0.2165 & 12.57  & 24  & 1866  & 19 & 37777  & W\\
$(1/4)^3$  & axis      & 32 & 0.0541 & 50.27  & 18  & 1084  & 15 & 25626  & W\\
$(1/4)^3$  & axis      & 64 & 0.0271 & 100.5  & 16  & 874   & 14 & 22600  & W\\
$(1/4)^3$  & face diag & 2  & 0.6124 & 4.443  & 60  & 11240 & 40 & 106341 & W\\
$(1/4)^3$  & body diag & 2  & 0.5000 & 5.441  & 44  & 6082  & 34 & 78134  & W\\
\midrule
slender & long axis  & 2  & 0.7078 & 0.393 & 90 & 25190 & 47  & 154116  & O\\
slender & long axis  & 3  & 0.4719 & 0.589 & 42 & 5548  & 29  & 70375   & W\\
slender & long axis  & 4  & 0.3539 & 0.785 & 30 & 2870  & 23  & 48269   & W\\
slender & long axis  & 8  & 0.1769 & 1.571 & 18 & 1084  & 15  & 24113   & W\\
slender & long axis  & 32 & 0.0442 & 6.283 & 10 & 390   & 10  & 12450   & W\\
slender & face diag  & 2  & 0.5005 & 0.555 & 46 & 6640  & 34  & 76426   & W\\
slender & body diag  & 2  & 0.5000 & 0.556 & 46 & 6640  & 34  & 76426   & W\\
slender & short axis & 2  & 11.325 & 0.025 & -- & --    & --  & --      & K\\
slender & short axis & 4  & 5.6624 & 0.049 & -- & --    & --  & --      & K\\
slender & short axis & 8  & 2.8312 & 0.098 & -- & --    & 207 & 3489254 & K\\
slender & short axis & 16 & 1.4156 & 0.196 & -- & --    & 70  & 459214  & K\\
slender & short axis & 32 & 0.7078 & 0.393 & 90 & 25190 & 34  & 116833  & O\\
slender & short axis & 64 & 0.3539 & 0.785 & 30 & 2870  & 21  & 45780   & W\\
\bottomrule
\end{tabular}
\end{center}

\noindent
A dash means the bound provides no certificate at all: the term sequence of
\eqref{eq:bnd:rem} does not decrease, which for the short axis of the slender cell is not
a defect of the bound but the truth, since $\rho > 1$ there and the expansion diverges.
Three readings of the table are worth stating.

\begin{itemize}[nosep,leftmargin=1.4em]
\item \textbf{$L$ is set by $\rho$, not by $k\mathbf{R}$.} At a fixed offset in cells the
selected $L$ moves by at most $4$ from $\lambda/32$ to $\lambda/4$ even though
$|k\mathbf{R}|$ changes by a factor eight. Frequency enters the \emph{cost} only through
$n_{\max}$ and through the cancellation of the $l$-sum, not through $L$.
\item \textbf{The far plateau is cheap and flat.} Beyond about eight cells $L$ settles at
$10$--$16$ for the cubes and the whole-box evaluation costs under $2~\mu$s per offset.
\item \textbf{Only the whole box ever wins on cost}, at every offset where its bound
applies. Subdivision is not a cost optimisation; it is the repair for the offsets where
the whole-box bound needs an $L$ that \code{Float64} cannot deliver.
\end{itemize}

\subsection{The band edges the shipped selector produces}
\label{sec:join:ship}

The same quantities from Theorem~A and Theorem~B at $\mathrm{tol} = 10^{-14}$,
$L_{\max} = 56$, as \code{verify.jl} part (e) prints them
(\code{tables/e\_terms.txt}, $192$ rows; ``cost'' is the number of complex multiply--adds
the route issues, $-1$ means the bound is not met at $L \le 56$). \code{route} is the
method the library actually takes.

\begin{center}\footnotesize
\begin{longtable}{llrrrrrrrc}
\toprule
shape & dir & $n$ & $\rho$ & $|k\mathbf{R}|$ & $L_{\rm wb}$ & cost & $L_{\rm oct}$ & cost & route \\
\midrule
\endfirsthead
\toprule
shape & dir & $n$ & $\rho$ & $|k\mathbf{R}|$ & $L_{\rm wb}$ & cost & $L_{\rm oct}$ & cost & route \\
\midrule
\endhead
\bottomrule
\endfoot
\texttt{c32} & \texttt{axis} & $2$ & $0.866$ & $0.393$ & $-1$ & $-1$ & $50$ & $98028$ & (ii) \\
\texttt{c32} & \texttt{axis} & $3$ & $0.5774$ & $0.589$ & $44$ & $3491$ & $30$ & $43036$ & (i) \\
\texttt{c32} & \texttt{axis} & $4$ & $0.433$ & $0.785$ & $32$ & $1892$ & $24$ & $29848$ & (i) \\
\texttt{c32} & \texttt{axis} & $6$ & $0.2887$ & $1.178$ & $22$ & $931$ & $19$ & $20272$ & (i) \\
\texttt{c32} & \texttt{axis} & $8$ & $0.2165$ & $1.571$ & $18$ & $641$ & $16$ & $16184$ & (i) \\
\texttt{c32} & \texttt{axis} & $16$ & $0.1083$ & $3.142$ & $14$ & $405$ & $13$ & $10976$ & (i) \\
\texttt{c32} & \texttt{axis} & $32$ & $0.0541$ & $6.283$ & $12$ & $307$ & $11$ & $8064$ & (i) \\
\texttt{c32} & \texttt{axis} & $64$ & $0.0271$ & $12.566$ & $10$ & $223$ & $10$ & $6776$ & (i) \\
\texttt{c32} & \texttt{face} & $2$ & $0.6124$ & $0.555$ & $48$ & $4132$ & $36$ & $49182$ & (i) \\
\texttt{c32} & \texttt{face} & $3$ & $0.4082$ & $0.833$ & $30$ & $1673$ & $24$ & $27902$ & (i) \\
\texttt{c32} & \texttt{face} & $4$ & $0.3062$ & $1.111$ & $24$ & $1096$ & $20$ & $20818$ & (i) \\
\texttt{c32} & \texttt{face} & $6$ & $0.2041$ & $1.666$ & $18$ & $641$ & $16$ & $15722$ & (i) \\
\texttt{c32} & \texttt{face} & $8$ & $0.1531$ & $2.221$ & $16$ & $516$ & $14$ & $12600$ & (i) \\
\texttt{c32} & \texttt{face} & $16$ & $0.0765$ & $4.443$ & $12$ & $307$ & $12$ & $9464$ & (i) \\
\texttt{c32} & \texttt{face} & $32$ & $0.0383$ & $8.886$ & $10$ & $223$ & $11$ & $8064$ & (i) \\
\texttt{c32} & \texttt{face} & $64$ & $0.0191$ & $17.772$ & $10$ & $223$ & $10$ & $6776$ & (i) \\
\texttt{c32} & \texttt{body} & $2$ & $0.5$ & $0.68$ & $36$ & $2371$ & $31$ & $36848$ & (i) \\
\texttt{c32} & \texttt{body} & $3$ & $0.3333$ & $1.02$ & $24$ & $1096$ & $21$ & $23576$ & (i) \\
\texttt{c32} & \texttt{body} & $4$ & $0.25$ & $1.36$ & $20$ & $779$ & $18$ & $18158$ & (i) \\
\texttt{c32} & \texttt{body} & $6$ & $0.1667$ & $2.041$ & $16$ & $516$ & $15$ & $13468$ & (i) \\
\texttt{c32} & \texttt{body} & $8$ & $0.125$ & $2.721$ & $14$ & $405$ & $14$ & $11179$ & (i) \\
\texttt{c32} & \texttt{body} & $16$ & $0.0625$ & $5.441$ & $12$ & $307$ & $11$ & $8064$ & (i) \\
\texttt{c32} & \texttt{body} & $32$ & $0.0312$ & $10.883$ & $10$ & $223$ & $10$ & $6776$ & (i) \\
\texttt{c32} & \texttt{body} & $64$ & $0.0156$ & $21.766$ & $10$ & $223$ & $10$ & $6776$ & (i) \\
\texttt{c8} & \texttt{axis} & $2$ & $0.866$ & $1.571$ & $-1$ & $-1$ & $50$ & $99736$ & (ii) \\
\texttt{c8} & \texttt{axis} & $3$ & $0.5774$ & $2.356$ & $46$ & $3805$ & $31$ & $46172$ & (i) \\
\texttt{c8} & \texttt{axis} & $4$ & $0.433$ & $3.142$ & $32$ & $1892$ & $24$ & $31052$ & (i) \\
\texttt{c8} & \texttt{axis} & $6$ & $0.2887$ & $4.712$ & $24$ & $1096$ & $20$ & $22456$ & (i) \\
\texttt{c8} & \texttt{axis} & $8$ & $0.2165$ & $6.283$ & $20$ & $779$ & $17$ & $18144$ & (i) \\
\texttt{c8} & \texttt{axis} & $16$ & $0.1083$ & $12.566$ & $16$ & $516$ & $15$ & $13468$ & (i) \\
\texttt{c8} & \texttt{axis} & $32$ & $0.0541$ & $25.133$ & $14$ & $405$ & $13$ & $10976$ & (i) \\
\texttt{c8} & \texttt{axis} & $64$ & $0.0271$ & $50.265$ & $14$ & $405$ & $13$ & $10976$ & (i) \\
\texttt{c8} & \texttt{face} & $2$ & $0.6124$ & $2.221$ & $50$ & $4473$ & $36$ & $51464$ & (i) \\
\texttt{c8} & \texttt{face} & $3$ & $0.4082$ & $3.332$ & $30$ & $1673$ & $25$ & $30450$ & (i) \\
\texttt{c8} & \texttt{face} & $4$ & $0.3062$ & $4.443$ & $24$ & $1096$ & $21$ & $23030$ & (i) \\
\texttt{c8} & \texttt{face} & $6$ & $0.2041$ & $6.664$ & $20$ & $779$ & $17$ & $17654$ & (i) \\
\texttt{c8} & \texttt{face} & $8$ & $0.1531$ & $8.886$ & $18$ & $641$ & $16$ & $14798$ & (i) \\
\texttt{c8} & \texttt{face} & $16$ & $0.0765$ & $17.772$ & $16$ & $516$ & $14$ & $12600$ & (i) \\
\texttt{c8} & \texttt{face} & $32$ & $0.0383$ & $35.543$ & $14$ & $405$ & $13$ & $10976$ & (i) \\
\texttt{c8} & \texttt{face} & $64$ & $0.0191$ & $71.086$ & $14$ & $405$ & $13$ & $10976$ & (i) \\
\texttt{c8} & \texttt{body} & $2$ & $0.5$ & $2.721$ & $38$ & $2631$ & $31$ & $38276$ & (i) \\
\texttt{c8} & \texttt{body} & $3$ & $0.3333$ & $4.081$ & $26$ & $1275$ & $22$ & $24164$ & (i) \\
\texttt{c8} & \texttt{body} & $4$ & $0.25$ & $5.441$ & $22$ & $931$ & $19$ & $20230$ & (i) \\
\texttt{c8} & \texttt{body} & $6$ & $0.1667$ & $8.162$ & $18$ & $641$ & $16$ & $15953$ & (i) \\
\texttt{c8} & \texttt{body} & $8$ & $0.125$ & $10.883$ & $16$ & $516$ & $15$ & $14336$ & (i) \\
\texttt{c8} & \texttt{body} & $16$ & $0.0625$ & $21.766$ & $14$ & $405$ & $14$ & $11179$ & (i) \\
\texttt{c8} & \texttt{body} & $32$ & $0.0312$ & $43.531$ & $14$ & $405$ & $13$ & $10976$ & (i) \\
\texttt{c8} & \texttt{body} & $64$ & $0.0156$ & $87.062$ & $14$ & $405$ & $13$ & $10976$ & (i) \\
\texttt{c4} & \texttt{axis} & $2$ & $0.866$ & $3.142$ & $-1$ & $-1$ & $51$ & $104384$ & (ii) \\
\texttt{c4} & \texttt{axis} & $3$ & $0.5774$ & $4.712$ & $48$ & $4132$ & $32$ & $49420$ & (i) \\
\texttt{c4} & \texttt{axis} & $4$ & $0.433$ & $6.283$ & $34$ & $2125$ & $26$ & $35224$ & (i) \\
\texttt{c4} & \texttt{axis} & $6$ & $0.2887$ & $9.425$ & $26$ & $1275$ & $21$ & $25900$ & (i) \\
\texttt{c4} & \texttt{axis} & $8$ & $0.2165$ & $12.566$ & $22$ & $931$ & $19$ & $21308$ & (i) \\
\texttt{c4} & \texttt{axis} & $16$ & $0.1083$ & $25.133$ & $18$ & $641$ & $17$ & $17164$ & (i) \\
\texttt{c4} & \texttt{axis} & $32$ & $0.0541$ & $50.265$ & $18$ & $641$ & $16$ & $15260$ & (i) \\
\texttt{c4} & \texttt{axis} & $64$ & $0.0271$ & $100.531$ & $16$ & $516$ & $15$ & $14336$ & (i) \\
\texttt{c4} & \texttt{face} & $2$ & $0.6124$ & $4.443$ & $52$ & $4827$ & $37$ & $54880$ & (i) \\
\texttt{c4} & \texttt{face} & $3$ & $0.4082$ & $6.664$ & $32$ & $1892$ & $26$ & $33740$ & (i) \\
\texttt{c4} & \texttt{face} & $4$ & $0.3062$ & $8.886$ & $26$ & $1275$ & $22$ & $27132$ & (i) \\
\texttt{c4} & \texttt{face} & $6$ & $0.2041$ & $13.329$ & $22$ & $931$ & $19$ & $20762$ & (i) \\
\texttt{c4} & \texttt{face} & $8$ & $0.1531$ & $17.772$ & $20$ & $779$ & $18$ & $18662$ & (i) \\
\texttt{c4} & \texttt{face} & $16$ & $0.0765$ & $35.543$ & $18$ & $641$ & $16$ & $16184$ & (i) \\
\texttt{c4} & \texttt{face} & $32$ & $0.0383$ & $71.086$ & $18$ & $641$ & $15$ & $14336$ & (i) \\
\texttt{c4} & \texttt{face} & $64$ & $0.0191$ & $142.172$ & $16$ & $516$ & $15$ & $14336$ & (i) \\
\texttt{c4} & \texttt{body} & $2$ & $0.5$ & $5.441$ & $40$ & $2904$ & $32$ & $41244$ & (i) \\
\texttt{c4} & \texttt{body} & $3$ & $0.3333$ & $8.162$ & $28$ & $1467$ & $23$ & $28392$ & (i) \\
\texttt{c4} & \texttt{body} & $4$ & $0.25$ & $10.883$ & $24$ & $1096$ & $20$ & $23548$ & (i) \\
\texttt{c4} & \texttt{body} & $6$ & $0.1667$ & $16.324$ & $20$ & $779$ & $18$ & $19180$ & (i) \\
\texttt{c4} & \texttt{body} & $8$ & $0.125$ & $21.766$ & $20$ & $779$ & $17$ & $18144$ & (i) \\
\texttt{c4} & \texttt{body} & $16$ & $0.0625$ & $43.531$ & $18$ & $641$ & $16$ & $16184$ & (i) \\
\texttt{c4} & \texttt{body} & $32$ & $0.0312$ & $87.062$ & $16$ & $516$ & $15$ & $14336$ & (i) \\
\texttt{c4} & \texttt{body} & $64$ & $0.0156$ & $174.125$ & $16$ & $516$ & $15$ & $14336$ & (i) \\
\texttt{sl} & \texttt{ax1} & $2$ & $0.7078$ & $0.393$ & $-1$ & $-1$ & $44$ & $78652$ & (ii) \\
\texttt{sl} & \texttt{ax1} & $3$ & $0.4719$ & $0.589$ & $38$ & $2631$ & $28$ & $38360$ & (i) \\
\texttt{sl} & \texttt{ax1} & $4$ & $0.3539$ & $0.785$ & $30$ & $1673$ & $23$ & $28476$ & (i) \\
\texttt{sl} & \texttt{ax1} & $6$ & $0.2359$ & $1.178$ & $22$ & $931$ & $18$ & $19180$ & (i) \\
\texttt{sl} & \texttt{ax1} & $8$ & $0.1769$ & $1.571$ & $18$ & $641$ & $16$ & $15260$ & (i) \\
\texttt{sl} & \texttt{ax1} & $16$ & $0.0885$ & $3.142$ & $14$ & $405$ & $13$ & $10976$ & (i) \\
\texttt{sl} & \texttt{ax1} & $32$ & $0.0442$ & $6.283$ & $12$ & $307$ & $11$ & $8064$ & (i) \\
\texttt{sl} & \texttt{ax1} & $64$ & $0.0221$ & $12.566$ & $10$ & $223$ & $10$ & $6776$ & (i) \\
\texttt{sl} & \texttt{ax2} & $2$ & $0.7078$ & $0.393$ & $-1$ & $-1$ & $44$ & $78652$ & (ii) \\
\texttt{sl} & \texttt{ax2} & $3$ & $0.4719$ & $0.589$ & $38$ & $2631$ & $28$ & $38360$ & (i) \\
\texttt{sl} & \texttt{ax2} & $4$ & $0.3539$ & $0.785$ & $30$ & $1673$ & $23$ & $28476$ & (i) \\
\texttt{sl} & \texttt{ax2} & $6$ & $0.2359$ & $1.178$ & $22$ & $931$ & $18$ & $19180$ & (i) \\
\texttt{sl} & \texttt{ax2} & $8$ & $0.1769$ & $1.571$ & $18$ & $641$ & $16$ & $15260$ & (i) \\
\texttt{sl} & \texttt{ax2} & $16$ & $0.0885$ & $3.142$ & $14$ & $405$ & $13$ & $10976$ & (i) \\
\texttt{sl} & \texttt{ax2} & $32$ & $0.0442$ & $6.283$ & $12$ & $307$ & $11$ & $8064$ & (i) \\
\texttt{sl} & \texttt{ax2} & $64$ & $0.0221$ & $12.566$ & $10$ & $223$ & $10$ & $6776$ & (i) \\
\texttt{sl} & \texttt{ax3} & $2$ & $11.3248$ & $0.025$ & $-1$ & $-1$ & $-1$ & $-1$ & (iii) \\
\texttt{sl} & \texttt{ax3} & $3$ & $7.5498$ & $0.037$ & $-1$ & $-1$ & $-1$ & $-1$ & (iii) \\
\texttt{sl} & \texttt{ax3} & $4$ & $5.6624$ & $0.049$ & $-1$ & $-1$ & $-1$ & $-1$ & (iii) \\
\texttt{sl} & \texttt{ax3} & $6$ & $3.7749$ & $0.074$ & $-1$ & $-1$ & $-1$ & $-1$ & (iii) \\
\texttt{sl} & \texttt{ax3} & $8$ & $2.8312$ & $0.098$ & $-1$ & $-1$ & $-1$ & $-1$ & (iii) \\
\texttt{sl} & \texttt{ax3} & $16$ & $1.4156$ & $0.196$ & $-1$ & $-1$ & $-1$ & $-1$ & (iii) \\
\texttt{sl} & \texttt{ax3} & $32$ & $0.7078$ & $0.393$ & $-1$ & $-1$ & $32$ & $59164$ & (ii) \\
\texttt{sl} & \texttt{ax3} & $64$ & $0.3539$ & $0.785$ & $30$ & $1673$ & $21$ & $27104$ & (i) \\
\end{longtable}

\end{center}

\noindent
Three things change against Section~\ref{sec:join:bands} and nothing else does. The
whole-box $L$ falls by two to six shells at two to four cells (Theorem~A is sharper than
\eqref{eq:bnd:rem} by a median factor $3.2\times10^{3}$) and rises by two shells in the far
plateau, where the old $L = 8$ was never a certificate of the sum the code performs; the
octant $L_j$ falls from $55$ to $50$ at $(2,0,0)$; and the cap of $56$ is reached by the
whole-box bound at the two-cell offsets, which is what routes them to (ii).

\paragraph{Routing counts over a whole block.} The number of offsets each route takes, over
the \code{egoToe} octant that Gila evaluates (indices with $\max_i i_i \ge 3$, i.e.\
max-norm separation $\ge 2$ cells), from \code{tables/e\_route.txt}:

\begin{center}\footnotesize
\begin{tabular}{llrrrrr}
\toprule
shape & $f$ & block & (i) & (ii) & (iii) & setup s \\
\midrule
\texttt{c32} & $1.0+0.0\ii$ & $32\!\times\!32\!\times\!32$ & $32748$ & $12$ & $0$ & $0.7$ \\
\texttt{c32} & $1.0+0.0\ii$ & $64\!\times\!64\!\times\!64$ & $262124$ & $12$ & $0$ & $0.7$ \\
\texttt{c32} & $1.0+0.0\ii$ & $128\!\times\!128\!\times\!128$ & $2097132$ & $12$ & $0$ & $0.8$ \\
\texttt{c32} & $1.0+0.1\ii$ & $32\!\times\!32\!\times\!32$ & $32748$ & $12$ & $0$ & $1.0$ \\
\texttt{c32} & $1.0+0.1\ii$ & $64\!\times\!64\!\times\!64$ & $262124$ & $12$ & $0$ & $0.9$ \\
\texttt{c32} & $1.0+0.1\ii$ & $128\!\times\!128\!\times\!128$ & $2097132$ & $12$ & $0$ & $0.9$ \\
\texttt{c8} & $1.0+0.0\ii$ & $32\!\times\!32\!\times\!32$ & $32748$ & $12$ & $0$ & $0.7$ \\
\texttt{c8} & $1.0+0.0\ii$ & $64\!\times\!64\!\times\!64$ & $262124$ & $12$ & $0$ & $0.8$ \\
\texttt{c8} & $1.0+0.0\ii$ & $128\!\times\!128\!\times\!128$ & $2097132$ & $12$ & $0$ & $0.8$ \\
\texttt{c8} & $1.0+0.1\ii$ & $32\!\times\!32\!\times\!32$ & $32748$ & $12$ & $0$ & $0.9$ \\
\texttt{c8} & $1.0+0.1\ii$ & $64\!\times\!64\!\times\!64$ & $262124$ & $12$ & $0$ & $0.9$ \\
\texttt{c8} & $1.0+0.1\ii$ & $128\!\times\!128\!\times\!128$ & $2097132$ & $12$ & $0$ & $1.1$ \\
\texttt{c4} & $1.0+0.0\ii$ & $32\!\times\!32\!\times\!32$ & $32748$ & $12$ & $0$ & $0.8$ \\
\texttt{c4} & $1.0+0.0\ii$ & $64\!\times\!64\!\times\!64$ & $262124$ & $12$ & $0$ & $0.8$ \\
\texttt{c4} & $1.0+0.0\ii$ & $128\!\times\!128\!\times\!128$ & $2097132$ & $12$ & $0$ & $0.9$ \\
\texttt{c4} & $1.0+0.1\ii$ & $32\!\times\!32\!\times\!32$ & $32748$ & $12$ & $0$ & $1.1$ \\
\texttt{c4} & $1.0+0.1\ii$ & $64\!\times\!64\!\times\!64$ & $262124$ & $12$ & $0$ & $1.0$ \\
\texttt{c4} & $1.0+0.1\ii$ & $128\!\times\!128\!\times\!128$ & $2097132$ & $12$ & $0$ & $1.1$ \\
\texttt{sl} & $1.0+0.0\ii$ & $16\!\times\!16\!\times\!32$ & $8013$ & $104$ & $67$ & $0.7$ \\
\texttt{sl} & $1.0+0.0\ii$ & $32\!\times\!32\!\times\!64$ & $65350$ & $111$ & $67$ & $0.7$ \\
\texttt{sl} & $1.0+0.0\ii$ & $64\!\times\!64\!\times\!128$ & $524102$ & $111$ & $67$ & $0.8$ \\
\texttt{sl} & $1.0+0.1\ii$ & $16\!\times\!16\!\times\!32$ & $8013$ & $104$ & $67$ & $0.9$ \\
\texttt{sl} & $1.0+0.1\ii$ & $32\!\times\!32\!\times\!64$ & $65350$ & $111$ & $67$ & $0.9$ \\
\texttt{sl} & $1.0+0.1\ii$ & $64\!\times\!64\!\times\!128$ & $524102$ & $111$ & $67$ & $0.9$ \\
\bottomrule
\end{tabular}

\end{center}

\noindent
\emph{On the three cubic shapes the $k$-series is never needed, and subdivision is needed at
exactly twelve offsets out of two million.} The whole-box $L$ histogram of the largest block
bottoms out at a flat floor --- $L = 10$ at $\lambda/32$, $14$ at $\lambda/8$, $16$ at
$\lambda/4$ --- which $97.8\,\%$, $99.6\,\%$ and $97.5\,\%$ of the offsets attain, and this
is what makes the average per-offset cost of a large build converge to a few hundred
nanoseconds.

\subsection{The twelve two-cell offsets of a cubic grid}
\label{sec:join:twelve}

On a cubic grid the offsets with $\max_i n_i \ge 2$ and $|\mathbf{n}| \le 2.56$ are
exactly twelve per octant, for every $N$:
\begin{gather*}
(2,0,0),\ (0,2,0),\ (0,0,2);\qquad
(2,1,1),\ (1,2,1),\ (1,1,2);\\
(2,1,0),\ (2,0,1),\ (1,2,0),\ (0,2,1),\ (1,0,2),\ (0,1,2).
\end{gather*}
These are precisely the offsets the rule sends to the octant split, at all three cubic
scales, and the reason is visible in the $L$ columns:

\begin{center}\footnotesize
\begin{tabular}{@{}lrrrl@{}}
\toprule
shape & $\mathbf{n}$ class & $\rho$ & $L_{\mathrm{wb}}$ & $L_j$ over the eight sub-boxes\\
\midrule
$(1/32)^3$ & $(2,0,0)$ & 0.8660 & 174 & $31,31,31,31,55,55,55,55$\\
$(1/32)^3$ & $(2,1,0)$ & 0.7746 & 104 & $28,28,31,31,38,38,55,55$\\
$(1/32)^3$ & $(2,1,1)$ & 0.7071 & 78  & $26,28,28,32,32,38,38,56$\\
$(1/8)^3$  & $(2,0,0)$ & 0.8660 & 176 & $32,32,32,32,55,55,55,55$\\
$(1/8)^3$  & $(2,1,0)$ & 0.7746 & 104 & $28,28,32,32,38,38,55,55$\\
$(1/8)^3$  & $(2,1,1)$ & 0.7071 & 78  & $26,29,29,32,32,39,39,56$\\
$(1/4)^3$  & $(2,0,0)$ & 0.8660 & 182 & $33,33,33,33,56,56,56,56$\\
$(1/4)^3$  & $(2,1,0)$ & 0.7746 & 108 & $30,30,33,33,39,39,56,56$\\
$(1/4)^3$  & $(2,1,1)$ & 0.7071 & 82  & $27,30,30,33,33,39,39,56$\\
\bottomrule
\end{tabular}
\end{center}

\noindent
Every one of the twelve has $L_{\mathrm{wb}} \ge 78 > L_{\max}$ and
$\max_j L_j \le 56 \le L_{\max}$, at all three cubic scales; the next offset outwards,
$(2,2,0)$ with $\rho = 0.6124$, has $L_{\mathrm{wb}} = 56$, $58$, $60$ at
$\lambda/32, \lambda/8, \lambda/4$ and is handled by the whole box. \emph{The rule
therefore reproduces, from the bounds alone, the split that was found empirically}:
octant for the twelve, whole box from $(2,2,0)$ outwards. Since twelve is independent of
$N$, subdividing them costs twelve tensor evaluations in a build of two million.

The shipped selector agrees, with room to spare: at $\mathrm{tol} = 10^{-14}$ Theorem~A
gives no certificate at $L \le 56$ for any of the twelve, and Theorem~B gives
$\max_j L_j = 50$, $50$ and $51$ at $(2,0,0)$ and $51$, $51$ and $52$ at $(2,1,1)$ for the
$\lambda/32$, $\lambda/8$ and $\lambda/4$ cubes --- below the cap $L_{\max} = 56$ in
every case, and delivering $4.4\times10^{-15}$ or better against the reference.
The routing counts of Section~\ref{sec:join:ship} are exactly twelve route-(ii) offsets per
block, at all three cubic scales, at $32^3$, $64^3$ and $128^3$ alike, and at both
frequencies.

\begin{remark}[a discrepancy in the sub-box $L$, and its size]
\label{rem:join:disc}
The earlier sub-box analysis records the bound as selecting $L_j \le 48$ at these
offsets, against $L_j \le 56$ here, both from the same written bound
\eqref{eq:bnd:sub}, the same geometry ($\rho_j = 0.522$ at $(2,0,0)$ in both) and the
same budget $\mathrm{tol}\cdot\mathrm{est}/8$. The difference is a factor of about $100$
in the value of the tail, i.e.\ $\rho_j^{\,7}$; it does not change any band edge, since
both numbers are below $L_{\max} = 64$ and both are above the measured requirement of
$L_{13} = 38$--$42$. I have not been able to reproduce the smaller value from the bound
as written, and I record the larger one because it is the one this document's script
computes from the stated formula. The corresponding over-selection factor against the
measured requirement is $56/42 = 1.33$ here and $48/42 = 1.14$ there.
\end{remark}

\subsection{The slender cell: which offsets the expansion cannot reach}
\label{sec:join:slender}

For the slender cell $(1/32, 1/32, 1/512)$ the isotropic ratio is
$\rho = r_{\mathrm{d}}/|\mathbf{R}|$ with $r_{\mathrm{d}} = 0.04421$, and along the short
axis $|\mathbf{R}| = n/512$, so $\rho = 22.63/n$: the expansion \emph{diverges} until
$n = 23$ and reaches $\rho < 0.5$ only at $n = 46$. This is not slow convergence, and no
constant repairs it: the difference box still spans $\pm1/32$ in $x$ and $y$ while the
offset advances in steps of $1/512$. Subdivision helps but not enough --- the octant
bound needs $L = 207$ at $(0,0,8)$ and $L = 70$ at $(0,0,16)$.

Enumerating every offset with $n_1, n_2 \le 2$ and $n_3 \le 64$ that is not a touching
offset --- $577$ of them --- and asking of each whether the whole-box bound
\eqref{eq:bnd:rem} \emph{or} the octant bound \eqref{eq:bnd:sub} reaches
$10^{-13}\mathrm{est}$ within a given cap:

\begin{center}\small
\begin{tabular}{@{}lrl@{}}
\toprule
cap & count & the set\\
\midrule
$L > 48$ & \textbf{83} &
$(0,0,n)$ for $2 \le n \le 21$;\ $(0,1,n)$, $(1,0,n)$, $(1,1,n)$ for $2 \le n \le 22$\\
$L > 64$ & \textbf{64} &
$(0,0,n)$, $(0,1,n)$, $(1,0,n)$, $(1,1,n)$ for $2 \le n \le 17$\\
\bottomrule
\end{tabular}
\end{center}

\noindent
Both sets are stable under the reflections and under $N$. Every offset with $n_1 = 2$ or
$n_2 = 2$ is already covered by the whole box or the octant split, and so is every
$(n_1,n_2,n_3)$ with $n_3$ beyond the band edge. \emph{This is the only band, on any of the
four shapes, where neither form of the expansion has a certificate.} The shipped selector,
with Theorem~A/B at $\mathrm{tol} = 10^{-14}$ and $L_{\max} = 56$, hands
\textbf{$67$ offsets} to the $k$-series --- the four classes $(0,0,n_3)$, $(0,1,n_3)$,
$(1,0,n_3)$, $(1,1,n_3)$ with $n_3$ up to $18$ --- and $111$ to the octant split, out of
$524\,280$ offsets of a $64\times64\times128$ block; the counts do not change between
$32\times32\times64$ and $64\times64\times128$, nor between $f = 1$ and $f = 1+0.1\ii$.

The $k$-series owns it comfortably, because the cell is short in the only direction that
matters to $k$. Recomputed on exactly the $67$ offsets the router sends there
(\code{verify.jl} part (f), \code{tables/f\_bands.txt}), with $D_{\max}$ the largest
$|\mathbf{x}-\mathbf{y}|$ over the thirty-six face pairs, $x = |k|D_{\max}$,
$\Lambda = \bigl(\sum_{\mathrm{fp}}I^{(\mathrm{fp})}_{-1}\bigr)/
(4\pi|f|^2V_{\mathrm{t}}\,\mathrm{est}(\mathbf{R}))$ the assembly amplification, and $N$
the order Theorem~\ref{thm:bnd:kser} needs for $10^{-16}/\Lambda$:
\[
x \in [0.280,\ 0.602], \qquad N \in [12,\ 17], \qquad
\Lambda \le 1002, \qquad \varepsilon\Lambda \le 2.23\times10^{-13},
\]
identically at $f = 1$ and $f = 1+0.1\ii$. \emph{That last number is why route (iii) is
done in $128$-bit arithmetic}: a \code{Float64} face-pair $k$-series on this set would be
$2.2\times10^{-13}$ wrong from rounding alone, before any truncation. The three cubic
shapes have no route-(iii) offsets at all, and the table says so explicitly. The price is
the cost: $1.9$ to $9.5$~s per offset on eight threads, $128$~s for all $67$ from cold, once
per (shape, frequency), cached thereafter (\code{ksrcache/}), against $0.14$~s for the
other $524\,213$ offsets of the block.

\subsection{The geometry that makes this affordable}
\label{sec:join:geom}

Two facts about the lattice decide the cost of the whole scheme.

\paragraph{The convergence ratio is scale free and tabulated.} For cubic cells
$\rho = \sqrt{3}\,s/|\mathbf{R}|$ depends only on $\mathbf{n}$, and
$\rho_m$, the ratio after splitting each axis into $m$ pieces, is
$\min_j |\mathbf{h}|/|\mathbf{R}+\mathbf{c}_j|$:

\begin{center}\footnotesize
\begin{tabular}{@{}lrrrrr@{}}
\toprule
class & $n$ & $\rho$ & $\rho_2$ & $\rho_4$ & $\rho_8$\\
\midrule
axis & 2 & 0.866025 & 0.522233 & 0.333333 & 0.190117\\
axis & 3 & 0.577350 & 0.333333 & 0.190117 & 0.101535\\
axis & 4 & 0.433013 & 0.242536 & 0.132453 & 0.069171\\
axis & 8 & 0.216506 & 0.114960 & 0.059655 & 0.030378\\
axis & 16 & 0.108253 & 0.055815 & 0.028387 & 0.014313\\
axis & 32 & 0.054127 & 0.027486 & 0.013856 & 0.006956\\
axis & 64 & 0.027063 & 0.013637 & 0.006846 & 0.003430\\
face diag & 2 & 0.612372 & 0.397360 & 0.242536 & 0.135665\\
face diag & 8 & 0.153093 & 0.081559 & 0.042220 & 0.021485\\
body diag & 2 & 0.500000 & 0.333333 & 0.200000 & 0.111111\\
body diag & 8 & 0.125000 & 0.066667 & 0.034483 & 0.017544\\
$(n,1,0)$ & 2 & 0.774597 & 0.522233 & 0.333333 & 0.190117\\
$(n,2,1)$ & 2 & 0.577350 & 0.397360 & 0.242536 & 0.135665\\
\bottomrule
\end{tabular}
\end{center}

\noindent
Refinement divides $\rho$ by roughly $m(n-1)+1$ rather than by $m$, which is why the
octant split ($m = 2$) is worth its factor of about $3.6$ in cost at two cells and nothing
at all beyond three.

\paragraph{The slow band has a fixed size.} The offsets with large $\rho$ form a ball of
fixed radius in $\mathbf{n}$, $n_1^2+n_2^2+n_3^2 < 3/\rho^2$, whose population does not
grow with the volume: in the octant of offsets Gila actually evaluates there are $37$
offsets with $\rho > 0.5$, $141$ with $\rho > 0.3$, $436$ with $\rho > 0.2$ and $3096$
with $\rho > 0.1$, \emph{identically for $N = 32$, $64$ and $128$}. So the expensive
offsets are an $O(1)$ set and the average per-offset cost of a large build converges to
the far-plateau cost of the table in Section~\ref{sec:join:bands}, a few hundred
nanoseconds. The slender cell is the same statement with a different constant: its
$k$-series band contains $67$ offsets whatever $N$ is, and its octant band $111$.
